\documentclass[10pt,a4paper]{amsart}
\usepackage[margin=19mm]{geometry}
\usepackage{amsmath,amssymb,mathtools}
\usepackage[scr=rsfs,cal=euler]{mathalfa}
\usepackage{xfrac}
\usepackage[unicode,hidelinks,bookmarks=false]{hyperref}
\newcommand{\ie}{i.e.\ }
\newcommand{\cf}{cf.\ }
\newcommand{\resp}{respectively\ }
\newcommand{\Subsection}{\subsection}
\providecommand{\nobreakdash}{\nobreak\hbox{-}}
\setlength{\emergencystretch}{2em}
\multlinegap0pt
\usepackage{bm}
\usepackage{bbm}
\usepackage{dsfont}
\usepackage{xspace}
\usepackage{cancel}
\usepackage{mathtools}
\usepackage{amsthm}
\makeatletter
\@ifundefined{thm}{\newtheorem{thm}{Thm}}{}
\@ifundefined{lem}{\newtheorem{lem}[thm]{Lemma}}{}
\makeatother
\def\wdt{\widetilde}
\newcommand{\bfu}{{\mathbf u}}
\newcommand{\bfv}{{\mathbf v}}
\newcommand{\bfx}{{\mathbf x}}
\newcommand{\mB}{{\mathcal B}}
\newcommand{\upM}{{\overline{M}}}
\newcommand{\ups}{{\overline{s} }}
\title[Box dimension of the graphs of recurrent fractal interpolati]{Box dimension of the graphs of recurrent fractal interpolation functions}
\author{Lai Jiang, Xiao-Hui Li, Zhen Liang, Huo-Jun Ruan}
\date{Source: arXiv:2510.02754v1 (CC BY 4.0)}
\begin{document}
\maketitle

\noindent\emph{Source and licence.} Box dimension of the graphs of recurrent fractal interpolation functions. Version arXiv:2510.02754v1 (CC BY 4.0), distributed under the Creative Commons Attribution 4.0 International licence, as recorded in the arXiv licence metadata for this exact version and shown on the abstract page https://arxiv.org/abs/2510.02754v1. Full rights notice: licenses/arxiv-2510.02754-CC-BY-4.0.txt.\par\smallskip

\noindent\textit{Editorial scope.} Counted unit: the Lem whose source label is \texttt{lem:3-rho-key-lem-1} in the article (source0.tex lines 746--750, source label \texttt{lem:3-rho-key-lem-1}), delivered together with the article's own complete original proof of it (source0.tex lines 752--844, one \texttt{proof} environment of 93 lines). No mathematics is written, repaired or re-derived: statement and proof are the article's own text line for line. Required context retained uncounted: lem:Theta-simple-facts (lines 497--510); lem:PFN (lines 654--659). Every definition, display and label of those lines is retained unchanged. Omitted from this package and not redistributed: 1--496, 511--653, 660--745, 751--751, 845--2092. Nothing inside a retained excerpt is omitted or altered: every retained body is delivered exactly as the article's lines are written, so no omission receipt of any kind accompanies this package. Citations: the retained excerpts carry no citation command at all, so no bibliography entry of the article is transcribed and no reference is redistributed. Reference adaptation: none was needed and none was made; every reference inside the delivered unit resolves inside it, so no unresolved reference or citation is produced. Printed slips kept literal: no printed word, symbol, hypothesis or number of the counted statement or of its proof was altered. Layout and class: the article's own class is \texttt{amsart} with packages including amsfonts, amssymb, amsmath, graphicx, psfrag, color, verbatim, epstopdf, geometry; the wrapper is the lane's pinned \texttt{amsart} editorial wrapper at 10pt on A4, which restates the theorem-like environments the retained text needs and adds the article's own macro definitions that the retained text needs (\texttt{\textbackslash bfu}, \texttt{\textbackslash bfv}, \texttt{\textbackslash bfx}, \texttt{\textbackslash mB}, \texttt{\textbackslash upM}, \texttt{\textbackslash ups}, \texttt{\textbackslash wdt}), with extra wrapper packages (bm, bbm, dsfont, xspace, cancel, mathtools). The delivered numbering is the unit's own. Assets: no retained span contains a figure environment, a graphics-inclusion command, a PDF-page inclusion, a table or a TikZ picture; no image file is copied into this package, the package declares no asset dependency and no image file is redistributed with it. Licence: version arXiv:2510.02754v1 of this article is distributed under the Creative Commons Attribution 4.0 International licence, as recorded in the arXiv licence metadata for this exact version and shown on the abstract page https://arxiv.org/abs/2510.02754v1. A full rights notice is delivered at licenses/arxiv-2510.02754-CC-BY-4.0.txt. Characters: every retained excerpt is pure ASCII, so the delivered engine, XeLaTeX with its default Latin Modern text font, reports no missing character.
\begin{lem}\label{lem:Theta-simple-facts}
 Let $1\leq r\leq m$ and $k\geq 1$. 
 \begin{enumerate}
	  \item For any two distinct elements $i,j$ in $\{1,2,\ldots,d_r{T_r}^{k-1}\}$, we have $(I_{r,i}^k \cap I_{r,j}^k)^\circ =\emptyset$, where we use $(\cdot)^\circ$ to denote the interior of a set.
	  \item If $i\in\{1,2,\ldots,d_r{T_r}^{k-1}\}\setminus \Theta_{r,k}$, then $(I_{r,i}^k \cap B_{r,k})^\circ = \emptyset$. As a result, $(I_{r,i}^k \cap B_{r,k+1})^\circ = \emptyset$.
	  \item $\Theta_{r,k}=\big\{ 1 \leq  i \leq d_r {T_r}^{k-1}: \,I_{r,i}^k \subset B_{r,k}  \big\}$.	 In particular, $\Theta_{r,1}=\{1,2,\ldots,d_r\}$.
	  \item For any $i\in \Theta_{r,k}$ and $k'>k$, the set $\{i'\in \Theta_{r,k'}:\, I_{r,i'}^{k'}\subset I_{r,i}^k\}$ is nonempty.
	  \item For any $1\leq t\leq d_r$, $D_{r,t}^1=D_{a_{r,t}}$. Furthermore, for any $1\leq i\leq d_r{T_r}^k$, there exist $n_i\in\{1,\ldots,N\}$ and $p_i\in\{1,2,\ldots,{T_r}^{k-1}\}$ dependent of $r$ and $k$, such that
	  \begin{equation}\label{eq:Theta-simple-facts-1}
	    D_{r,i}^{k+1}=[x_{(n_i-1)}+(p_i-1)\varepsilon_{r,k},x_{(n_i-1)}+p_i\varepsilon_{r,k}],
	  \end{equation}
	 where $\varepsilon_{r,k}=(x_N-x_0)N^{-1}{T_r}^{-k+1}$. As a result, we have either $D_{r,i}^{k+1}\in \mB_{r,k}$ or $(D_{r,i}^{k+1}\cap B_{r,k})^\circ=\emptyset$.
  \end{enumerate}	
\end{lem}

\begin{lem}\label{lem:PFN}
Let $A$ be a nonnegative matrix. 
Then $\rho(A)$ is an eigenvalue of $A$ and there
is a nonnegative nonzero vector $\bfx$ such that $A\bfx  = \rho(A)\bfx$.

\end{lem}

\begin{lem}\label{lem:3-rho-key-lem-1}
For any $1\leq r\leq m$, $k \geq 1$ and $\Omega \subset \{ 1, \ldots, d_r {T_r}^{k-1} \}$, we
have
$$ \rho(\upM_{r,k}|_{\Omega}) = \rho( \upM_{r,k}^*|_{\wdt\Omega}).  $$
\end{lem}

\begin{proof}
%For any $i \in \Omega$, we define 
%$$\I_{r,k}(i):=\{j\in \Omega: I_{r,j}^k \subset D_{r,i}^k \}.$$
%Assume that $D_i^k=I_\ell^{k-1}$ where $1 \leq \ell \leq NT^{k-2}$, we have 
%$$\I_k(i)=\{(\ell-1)T+1, \ldots, \ell T \}.$$

Firstly, we are going to prove $\rho(\upM_{r,k} |_\Omega)\geq \rho(\upM_{r,k}^* |_{\wdt{\Omega}})$.
Write $\lambda=\rho(\upM_{r,k}^*|_{\wdt{\Omega}} )$.
From Lemma~\ref{lem:PFN}, 
$\lambda$ is an eigenvalue of $\upM_{r,k}^* |_{\wdt{\Omega}}$ and there is a nonnegative nonzero vector 
$ \wdt\bfu=(\wdt{u}_{i'})_{i'\in \wdt{\Omega}}$ such that $(\upM_{r,k}^* |_{\wdt{\Omega}}) \wdt\bfu^T =\lambda \wdt\bfu^T$.

Define a vector $\bfu=(u_i)_{i\in\Omega}$ by letting
\[
{u}_i=\sum_{i'=(i-1)T_r+1}^{iT_r} \wdt{u}_{i'}, \quad i\in\Omega.
\]
Obviously, $\bfu$ is also nonnegative and nonzero. 


Now we fix $i\in \Omega$. For any $(i-1)T_r<i'\leq iT_r$, by the definition of $\wdt\bfu^T$,
\[
  \lambda \wdt{u}_{i'} =\sum_{j'\in \wdt{\Omega}} \big(\upM_{r,k}^*\big)_{i',j'}  \wdt{u}_{j'}.
\]
Let $n$ be the unique element in $\{1,2,\ldots,N\}$ satisfying $I_{r,i'}^{k+1}\subset I_{r,i}^k\subset I_n$.

In the case that $D_{r,i'}^{k+1}\not\subset \bigcup_{p\in \Lambda_r} I_p$, it is clear that for all $j'\in\wdt{\Omega}$, we have $I^{k+1}_{r,j'}\not\subset D_{r,i'}^{k+1}$ so that $\big(\upM_{r,k}^*\big)_{i',j'}=0$. Hence, $\lambda \wdt{u}_{i'}=0$.  

In the case that $D_{r,i'}^{k+1}\subset \bigcup_{p\in \Lambda_r} I_p$, from Lemma~\ref{lem:Theta-simple-facts}(5), there exists $\ell\in\{1,\ldots,d_r{T_r}^{k-1}\}$ such that $I_{r,\ell}^k=D_{r,i'}^{k+1}$. If $\ell\not\in\Omega$, then for all $j'\in\wdt{\Omega}$, we have $I^{k+1}_{r,j'}\not\subset I_{r,\ell}^k= D_{r,i'}^{k+1}$ so that $\big(\upM_{r,k}^*\big)_{i',j'}=0$, which gives $\lambda \wdt{u}_{i'}=0$. If $\ell\in\Omega$, then 
$\big(\upM_{r,k}^*\big)_{i',j'}\not=0$ if and only if $I^{k+1}_{r,j'}\subset D_{r,i'}^{k+1}=I^k_{r,\ell}$, which is equivalent to $(\ell-1)T_r<j'\leq \ell T_r$. Furthermore, if $(\ell-1)T_r<j'\leq \ell T_r$, then
\[
\big(\upM_{r,k}^*\big)_{i',j'}=\max \big\{|S_n(x)|:\, x\in D^{k+1}_{r,i'} \big\}
=\max\big\{|S_n(x)|:\, x\in I^k_{r,\ell}\big\}=\ups_{r,i,\ell}^k,
\]
where the last equality follows from the fact that $I_{r,\ell}^k=D_{r,i'}^{k+1}\subset D_{r,i}^k$. Hence
\[
\lambda \wdt{u}_{i'} = \ups_{r,i,\ell}^k \sum_{j'=(\ell-1)T_r+1}^{ \ell T_r } \wdt{u}_{j'}
=\ups_{r,i,\ell}^k  {u}_{\ell}.
\]

%$\I_{r,k}(i)
%\begin{align*}
%	\{i':\, (i-1)T_r<i'\leq iT_r \mbox{ and } D_{r,i'}^{k+1}\subset \bigcup_{p\in \Lambda_r} I_p\}  &\to \{\ell\in \Omega: I_{r,\ell}^k \subset D_{r,i}^k \} \\
%	i' &\mapsto \ell
%\end{align*}

Notice that the map
$i'\mapsto \ell$ defined by $I^{k}_{r,\ell}=D^{k+1}_{r,i'}$ is a bijection from 
\[
\big\{i':\, (i-1)T_r<i'\leq iT_r \mbox{ and } D_{r,i'}^{k+1}=I_{r,\ell}^k \mbox{ for some } \ell\in \Omega\big\}
\] 
to $\{\ell\in \Omega: I_{r,\ell}^k \subset D_{r,i}^k \}$. Thus from the above arguments, 
\[
  \lambda {u}_i=\sum_{i'=(i-1)T_r+1}^{iT_r} \lambda \wdt{u}_{i'}=\sum_{\ell\in \Omega: I_{r,\ell}^k \subset D_{r,i}^k} \ups_{r,i,\ell}^k u_{\ell}=
   \sum_{\ell\in \Omega} (\upM_{r,k})_{i,\ell}  u_{\ell},
\]
where the last equality follows from the definitions of $\upM_{r,k}$ and $\ups_{r,i,\ell}^k$.
%\big((\upM_{r,k}|_\Omega )\bfu^T \big)_{i}.


From the arbitrariness of $i$, we have $\lambda \bfu^T=(\upM_{r,k} |_\Omega) \bfu^T$.
Thus, $\lambda$ is an eigenvalue of $\upM_{r,k}|_\Omega $, which implies $\rho(\upM_{r,k}^* |_{\wdt{\Omega}} )=\lambda \leq \rho(\upM_{r,k}|_\Omega)$.


In the second part, we are going to prove $\rho(\upM_{r,k}|_\Omega)\leq \rho(\upM_{r,k}^* |_{\wdt{\Omega}})$. Without loss of generality, we assume that 
$\mu=\rho(\upM_{r,k}|_\Omega)>0$. 
From Lemma~\ref{lem:PFN}, there is a nonnegative nonzero vector $\bfv=(v_i)_{i\in\Omega}$ such that $\big( \upM_{r,k} |_\Omega \big)\bfv^T =\mu \bfv^T$.

For any $i\in \Omega$ and $(i-1)T_r<i'\leq iT_r$, if there exists $\ell \in \Omega$ such that $I_{r,\ell}^k=D_{r,i'}^{k+1}$, then we define $\wdt{v}_{i'}=	\ups_{r,i,\ell}^k v_{\ell}$; otherwise, we define $\wdt{v}_{i'}=0$. We remark that in the former case, $I_{r,\ell}^k\subset D_{r,i}^k$. Hence,
it is straightforward to see that
\begin{equation}\label{eq:3-eigen-vv-sum}
	\sum_{i'=(i-1)T_r+1}^{iT_r} \wdt{v}_{i'}=\sum_{\ell\in \Omega: I_{r,\ell}^k \subset D_{r,i}^k} \ups_{r,i,\ell}^k v_{\ell}=\sum_{\ell\in \Omega} (\upM_{r,k})_{i,\ell}  v_{\ell}=\mu v_i.
\end{equation}
As a result, the vector $\wdt{\bfv}=(\wdt{v}_{i'})_{i'\in\wdt{\Omega}}$ is nonnegative and nonzero.


Fix $i\in\Omega$ and $i'\in\{(i-1)T_r+1,\ldots,iT_r\}$. Assume that $D_{r,i'}^{k+1}\not\subset \bigcup_{p\in \Lambda_r} I_p$ or $D_{r,i'}^{k+1}=I_{r,\ell}^k$ for some $\ell\not\in \Omega$. By definition, $\wdt{v}_{i'}=0$.
Notice that in this case, for all $j'\in \wdt{\Omega}$, we have $I^{k+1}_{r,j'}\not\subset D_{r,i'}^{k+1}$ so that $\big(\upM_{r,k}^*\big)_{i',j'}=0$. Hence, 
\[
\sum_{j'\in \wdt{\Omega}} \big(\upM_{r,k}^*\big)_{i',j'} \wdt{v}_{j'}=0=\mu \wdt{v}_{i'}.
\]

Assume that $D_{r,i'}^{k+1}=I_{r,\ell}^k$ for some $\ell\in \Omega$. By arguments in the first part of the proof, $\big(\upM_{r,k}^*\big)_{i',j'}\not=0$ if and only if $(\ell-1)T_r<j'\leq \ell T_r$, and if this holds, $\big(\upM_{r,k}^*\big)_{i',j'}=\ups_{r,i,\ell}^k$. Thus,
\[
\sum_{j'\in \wdt{\Omega}} \big(\upM_{r,k}^*\big)_{i',j'} \wdt{v}_{j'}=\sum_{j'=(\ell-1)T_r+1}^{\ell T_r} \big(\upM_{r,k}^*\big)_{i',j'} \wdt{v}_{j'}=\ups_{r,i,\ell}^k \sum_{j'=(\ell-1)T_r+1}^{ \ell T_r }  \wdt{v}_{j'}. 
\]
Combining this with \eqref{eq:3-eigen-vv-sum}, 
\[
\sum_{j'\in \wdt{\Omega}} \big(\upM_{r,k}^*\big)_{i',j'} \wdt{v}_{j'}=\ups_{r,i,\ell}^k \mu v_{\ell}=\mu \wdt{v}_{i'}.
\]

From the arbitrariness of $i'$, we have $\upM_{r,k}^*  \wdt{\bfv}^T =\mu \wdt{\bfv}^T$ so that $\mu$ is an eigenvalue of $\upM_{r,k}^*$. Hence, 
$\rho(\upM_{r,k})=\mu\leq \rho(\upM_{r,k}^*)$.
\end{proof}	

\end{document}
